\documentclass[11pt]{article}
\usepackage{fontspec}
\usepackage[french,english]{babel}
\usepackage{amsmath}
\usepackage{amssymb}
\usepackage{amsthm}
\usepackage{mathrsfs}
\usepackage{tikz-cd}
\usepackage{enumitem}

\newcommand{\printedpage}[1]{\marginpar{\footnotesize p.~#1}}
\newcommand{\msp}[1]{\marginpar{\footnotesize [#1]}}
\newcommand{\cercle}[1]{\tikz[baseline=(c.base)]\node[draw,circle,inner sep=0.8pt](c){#1};}
\newcommand{\Cat}{\mathrm{Cat}}
\newcommand{\id}{\mathrm{id}}
\newcommand{\inj}{\mathrm{inj}}
\newcommand{\Fib}{\mathrm{Fib}}
\newcommand{\Fl}{\mathrm{Fl}}

\begin{document}

\printedpage{64}
% head: A. Grothendieck : Dérivateurs, Ch. XV, Édition des Universités de Montpellier II et Paris VII
% foot: Version 5 novembre 2010

\msp{page 117}

\begin{description}[labelwidth=5.5em,leftmargin=6.5em,style=sameline]
\item[Loc\,(4)] (Cocartesian square axiom). Let $X$ be in Cat, $X_1$, $X_2$ two opens with union $X$, $X_0 = X_1\cap X_2$, whence a cocartesian diagram in Cat,
\[
\begin{tikzcd}
X_0 \arrow[r, "i_1"] \arrow[d, "i_2"'] & X_1 \arrow[d, "j_2"] \\
X_2 \arrow[r, "j_1"] & X\ .
\end{tikzcd}
\]
We then have $i_1\in W \iff j_1\in W$ (and symmetrically $i_2\in W \iff j_2\in W$).

\item[Loc\,(6)]\footnote[44]{There will be a more ``basic'' Loc\,(5) below.} (Axiom of $W$-fibrations). There are variants \cercle{a}, (a$'$), \cercle{b}, (b$'$), (b$''$), (c), (c$'$), \textit{cf.} XII, p. 53-59. The strongest form \cercle{a} is this:

Let $f : X\longrightarrow Y$ in Cat be a $W$-fibration in the old sense, \textit{i.e.} such that for every $Y''\longrightarrow Y'$ in $\Cat/Y$, with $Y'$, $Y''$ contractible, the arrow $X''\longrightarrow X'$ obtained by the base change $X\longrightarrow Y$ is in $W$. Then $f$ is in $\Fib_W$, \textit{i.e.} the base change functor by $f$, $\Cat/Y\longrightarrow\Cat/X$, transforms arrows in $W$ into arrows in $W$.

\item[Loc\,(7)] (Direct limit axiom). $W$ is stable under filtering direct limits.
\end{description}

Strengthened form of Loc\,(7):

\begin{description}[labelwidth=5.5em,leftmargin=6.5em,style=sameline]
\item[Loc\,(7\,bis)] (L-accessibility). $W$ is stable under filtering $\varinjlim$ \textit{and} accessible, \textit{i.e.} there exists a small subset $W_\pi$ of $\Fl(\Cat)$, such that $W$ is the smallest subset of $\Fl(\Cat)$ stable under filtering $\varinjlim$, containing $W_\pi$.
\end{description}

\msp{page 118}

\begin{description}[labelwidth=5.5em,leftmargin=6.5em,style=sameline]
\item[Loc\,(8)]\footnote[45]{I have changed the numbering, to take account of a Loc\,(5) below, which is again a stability axiom \texttt{[originally the axioms Loc\,(6) to Loc\,(8) were denoted Loc\,(5) to Loc\,(7)]}.} (Connectedness axiom). $f\in W\Longrightarrow \pi_0(f)$ bijective.
\end{description}

\textbf{N.B.} All the axioms, except the last (and a smidgen of the next to last) are stability axioms – they all say that $W$ is ``large'': it contains certain arrows, and if it contains such and such arrows, it contains such and such others. Loc\,(8), and to a certain extent Loc\,(7\,bis), are ``limiting'' axioms on the size of $W$: Loc\,(7) tells us that $W\subset\mathbf{W}_0$ (= set % [sic] Loc (7) -> Loc (8), apparently: the inclusion in W_0 is the axiome de connexité Loc (8)
of $f\in\Fl\,\Cat$ such that $\pi_0(f)$ is bijective), Loc\,(7\,bis) tells us of the existence of a cardinal such that $W$ is equal to the L-envelope of $W^\pi$ (hence $W\subset\text{L-env}(W^\pi)$).

I have tried to arrange the axioms in order of practical importance. I ended up adding the connectedness axiom at the end, for in the theoretical developments on general localizers, % [sic] théorique -> théoriques, agreement with "développements"
it practically never comes in. But I recall (XII, p. 20, th. 2) that if $W$ satisfies Loc\,(1, 2, 3, 6), and if Loc\,(8) is \textit{not} satisfied, then $W=\Fl(\Cat)$\,; thus in the presence of Loc\,(1, 2, 3, 6), Loc\,(8) \textit{is equivalent} to $W\neq\Fl(\Cat)$ in its entirety.

\msp{page 119}

The axioms that were used for 8.4, 1°, 3° and 8.5, $(b)$ are
\[
\mathrm{Loc}\,(1\,\mathrm{bis},\ 2,\ 3)\ .
\]

\printedpage{65}
% head: A. Grothendieck : Dérivateurs, Ch. XV, Édition des Universités de Montpellier II et Paris VII
% foot: Version 5 novembre 2010

Those which are necessary for 8.5, $(a)$ are
\[
\mathrm{Loc}\,(1\,\mathrm{bis},\ 2,\ 3,\ 4)\ .
\]
Finally, those necessary for 8.4, 2° are
\[
\mathrm{Loc}\,(1\,\mathrm{bis},\ 2,\ 3,\ 4,\ 7\,\mathrm{bis})\ ,
\]
and they imply the \textit{strong saturation of $W$} (\textit{cf.} next paragraph).

I must repair an omission, having somewhat treated with contempt (!) the questions of connected components, up to now.

\textbf{Proposition 9.1.} \textit{Suppose that $W$ satisfies} Loc\,(1, 2, 3, 4) \textit{and no more. Let $f_i : X_i\longrightarrow Y_i$ $(i\in I)$ be a} finite \textit{family of arrows of} Cat. \textit{For $f=\coprod f_i : \coprod X_i\longrightarrow\coprod Y_i$ to be in $W$, it suffices that $\forall\ i\in I$, we have $f_i\in W$. If $W$ also satisfies} Loc\,(7) (\textit{direct limit axiom}), \textit{then this conclusion remains valid without a finiteness condition on $I$.}

\textsc{Proof}. The case of finite $I$ is handled by induction on $n=\operatorname{card}I$, we are immediately reduced to the case where $n=2$\footnote[46]{The infinite case reduces to this under the hypothesis Loc\,(7), by passage to the limit.}. We first see the

\msp{page 120}

\textbf{Corollary 9.2.} \textit{Under the hypothesis} Loc\,(1, 4), \textit{let $f : X\longrightarrow Y$ be an arrow of} Cat, \textit{and $Z$ an object of} Cat. \textit{Then $f\in W$ if and only if $f\amalg\id_Z : X\amalg Z\longrightarrow Y\amalg Z$ is in $W$.}

Immediate from the cocartesian square, where $\inj_0^X$ is an open Dwyer immersion,
\[
\begin{tikzcd}[column sep=large, row sep=large]
X \arrow[r, "f"] \arrow[d, "\inj_0^X"'] & Y \arrow[d, "\inj_0^Y"] \\
X\amalg Z \arrow[r, "f'=f\amalg\id_Z"'] & Y\amalg Z\ ,
\end{tikzcd}
\]
which shows that $f\in W\iff f'\in W$.

By iterating this result, we find, if
\[
f_1\ :\ X_1\longrightarrow Y_1\ ,\qquad f_2\ :\ X_2\longrightarrow Y_2
\]
are given, that
\[
f_1\ \in\ W\ ,\ f_2\ \in\ W\quad\Longrightarrow\quad f_1\amalg f_2\ \in\ W\ ,
\]
by looking at the diagram of two cocartesian squares:
\[
\begin{tikzcd}[column sep=large, row sep=large]
X_1 \arrow[r, "f"] \arrow[d] & Y_1 \arrow[d] \\ % [sic] f -> f_1, the arrow X_1 -> Y_1 is f_1 (cf. f'_1 = f_1 ∐ id_{X_2})
X_1\amalg X_2 \arrow[r, "f'_1=f_1\amalg\id_{X_2}"'] \arrow[d, "f'_2=\id_{X_1}\amalg f_2"'] & Y_1\amalg X_2 \arrow[d, "g_2=\id_{Y_1}\amalg f_2"] \\
X_1\amalg Y_2 \arrow[r, "g_1=f_1\amalg\id_{Y_2}"'] & Y_1\amalg Y_2\ .
\end{tikzcd}
\]

\printedpage{66}
% head: A. Grothendieck : Dérivateurs, Ch. XV, Édition des Universités de Montpellier II et Paris VII
% foot: Version 5 novembre 2010

If $f_1$ and $f_2$ are in $W$, the same is true of $f'_1$, $f'_2$, $g_1$, $g_2$, by 9.2, hence also of $f_1\amalg f_2=g_2f'_1$.

\msp{page 121}

But I see no way of proving the converse, that $f_1\amalg f_2\in W\Longrightarrow f_1,f_2\in W$. Thanks to what has been proved, we are reduced to the case where $f_1$, $f_2$ are open Dwyer immersions, but it does not seem that this gets us a hair's breadth further\,! In short, we need to know that the canonical maps $Y_i\hookrightarrow\coprod Y_i$ are in $\Fib_W$. Since $Y_i$ is a direct summand, we are reduced to the case of a canonical immersion $Y_1\hookrightarrow Y_1\amalg Y_2=Y$: if $f : X\longrightarrow Y$ is in $W$, is the same true of $X_1=X\times_Y Y_1\longrightarrow Y_1$\,? This is a property which, of course, would be a consequence of the old W\,(4\,a) (XII, p. 53), or just W\,(4\,a$'$), hence of the current Loc\,(6\,a$'$). Even the form W\,(4\,c$'$) $=$ Loc\,(6\,c$'$), the weakest of all, suffices, as we see by factoring $X_1\longrightarrow Y$ as $X\xrightarrow{\ i\ }X'\xrightarrow{\ p\ }Y$, with $i\in W$ (more precisely, $i_1\ :\ X_1\longrightarrow X'_1$ and $i_2\ :\ X_2\longrightarrow X'_2$ in $W$\,!), and $p\ \in\ \mathrm{Parf}_W$. By Loc\,(6\,c$'$), $f\in W$, hence $p\in W$ implies that the fibers of $p$ are $W$-aspheric, which implies that $p$ is in $W_{\mathrm{univ}}$, hence $p_1 : X'_1\longrightarrow Y_1$ is in $W$.

\msp{page 122}

So let us note for conscience's sake:

\textbf{Corollary 9.3.} \textit{If $W$ satisfies} Loc\,(1, 2, 3, 6\,c$'$) (\textit{where} Loc\,(6\,c$'$) \textit{is the weakest of all the variants of the fibration axiom}), \textit{then for any family $(f_i : X_i\longrightarrow Y_i)_{i\in I}$ of arrows in} Cat, \textit{$f=\coprod f_i\in W\Longrightarrow$ the $f_i$ are in $W$ $\forall\ i\in I$}\footnote[47]{In fact, I do not know how to prove the implication in the reverse direction, assuming only Loc\,(1, 2, 3, 6\,c$'$).}. \textit{And if $i : Y_0\hookrightarrow Y$ is an arrow which is both an open immersion and a closed immersion} (i.e. \textit{an isomorphism with a direct summand of $Y$}), \textit{then $i$ is in} $\Fib_W$, i.e. \textit{if $f : X\longrightarrow X'$ is an arrow in} $\Cat/Y$, \textit{then $f\in W$ implies $(f_0 : X_0\longrightarrow Y_0)\in W$, where $X_0=X|Y_0$, $X'_0=X'|Y_0$.} % [sic] f_0 : X_0 -> Y_0 -> f_0 : X_0 -> X'_0, apparently: f is a flèche X -> X' and X'_0 is defined but otherwise unused

But it is rather scandalous to have to use such a strong axiom, to conclude from it such an innocuous result\,! I note this:

\textbf{Proposition 9.4.} \textit{Let $W\subset\Fl(\Cat)$, satisfying} Loc\,(1) \textit{and no more. The following conditions are equivalent} (\textit{and implied by} Loc\,(1, 2, 3, 4, 6\,c$'$)).
\begin{enumerate}[label={}, leftmargin=2.5em]
\item[a)] \textit{If $f_1 : X_1\longrightarrow Y_1$, $f_2 : X_2\longrightarrow Y_2$ are two arrows of} Cat, \textit{then}
\[
f_1\amalg f_2\in W\iff f_1\in W,\ f_2\in W\,.
\]
\item[b)] \textit{The following two conditions on $W$ are satisfied:}

\msp{page 123}

\begin{enumerate}[label={}, leftmargin=2.5em]
\item[1°)] \textit{If $f\ :\ X\longrightarrow Y$ is in $W$, then it is so couniversally under $\emptyset$,} i.e. \textit{for every $Z$ in} Cat, \textit{$f\amalg\id_Z : X\amalg Z\longrightarrow Y\amalg Z$ is in $W$.} (\textbf{N.B.} \textit{This is equivalent to: every finite sum of arrows $f_i\in W$ is in $W$ – and likewise even for infinite sums, if} Loc\,(7), \textit{the limit axiom, is satisfied.})
\item[2°)] \textit{Let $Y_0\hookrightarrow Y$ be an inclusion of a direct summand of an object of} Cat. \textit{Then if $X\xrightarrow{\ f\ }Y$ is in $W$, so is $X_0=X|Y_0\longrightarrow Y_0$.} (\textbf{N.B.} \textit{This is equivalent to the condition that the inclusion of a direct summand is always a $W$-fibration.})
\end{enumerate}

\textit{Moreover, if $\pi_0(Y)$ is finite, or if $W$ satisfies the limit axiom} Loc\,(7), \textit{it suffices to require 2° for $Y_0$ a connected component of $Y$.}

\printedpage{67}
% head: A. Grothendieck : Dérivateurs, Ch. XV, Édition des Universités de Montpellier II et Paris VII
% foot: Version 5 novembre 2010

\item[c)] \textit{For every} finite \textit{family $(f_i)_{i\in I}$ of arrows $f_i : X_i\longrightarrow Y_i$ in} Cat,
\[
\coprod f_i\in W\iff f_i\ \in\ W\ \ \textit{for all}\ i\in I.
\]
\end{enumerate}
\textit{Moreover, if $W$ satisfies} Loc\,(7), \textit{then it is unnecessary in} (c) \textit{to restrict ourselves to finite $I$.} (\textit{And in any case, the equivalent conditions} (a), (b), (c) \textit{imply for} every \textit{family of arrows $(f_i)_{i\in I}$, $\coprod f_i\in W\Longrightarrow f_i\in W$ $\forall\ i\in I$.})

\msp{page 124}

According to 9.3, these conditions are satisfied assuming Loc\,(1, 2, 3, 4, 6\,c$'$), but Loc\,(6\,c$'$) is a steam hammer for a fly\,! One would rather like to add an axiom to the axioms Loc\,(1) to Loc\,(4) and Loc\,(6) to Loc\,(8), which together with Loc\,(8) would be the most innocuous axiom of all.

\begin{description}[labelwidth=5.5em,leftmargin=6.5em,style=sameline]
\item[Loc\,(5)] (Decomposition axiom, weak form). If $(f_i)_{i\in I}$ is a finite family of arrows of Cat, then $f=\coprod f_i$ is in $W$ if and only if the $f_i$ are. (These are the equivalent conditions of 9.4.)\footnote[48]{It may also be convenient to introduce the still weaker condition:
\begin{description}[nosep,labelwidth=4.5em,leftmargin=5em,style=sameline]
\item[Loc\,(5\,a)] If $f_1$, $f_2$ are in $W$, then $f_1\amalg f_2\in W$.\\
This is also condition $(b)$, 1° of proposition 9.4.
\item[Loc\,(5\,b)] If $f_1\amalg f_2$ is in $W$, then $f_1$ and $f_2$ too.\\
$\iff$ If $Y_0\subset Y$ is a direct summand, then $Y_0\hookrightarrow Y$ is in $\Fib_W$.
\end{description}}

\item[Loc\,(5\,bis)] (Decomposition axiom, strong form). Like Loc\,(5), but without assuming $I$ finite.\footnote[49]{\texttt{[Let us rewrite:]}
\begin{description}[nosep,labelwidth=4.5em,leftmargin=5em,style=sameline]
\item[Loc\,(5)] \textit{a}) If $f_1,f_2\in W$, then $f_1\amalg f_2\in W$.\\
\textit{b}) If $f_1\amalg f_2\in W$, then $f_1,f_2\in W$\\
($\iff b'$) every direct summand $Y_0\hookrightarrow Y$ is in $\Fib_W$\,).
\item[Loc\,(5\,bis)] Like Loc\,(5), but $(a)$ replaced by ($a$\,bis) a family (not necessarily finite).
\end{description}}
\end{description}

Thus, assuming the limit axiom Loc\,(7), the two preceding axioms are equivalent. Moreover, let us note this:

\textbf{Corollary 9.5} (of proposition 9.4). \textit{The axiom} Loc\,(5\,bis) \textit{is equivalent to each of the following properties:}
\begin{enumerate}[label={}, leftmargin=3em]
\item[\cercle{a}] \textit{For every $Y$ of} Cat, \textit{decomposed into its connected components $Y_i$ $(i\in I=\pi_0(Y))$, and for every arrow $f\ :\ X\longrightarrow Y$ in} Cat, \textit{we have $f\ \in\ W$ if and only if the $f_i : X_i=X|Y_i\longrightarrow Y_i$ are in $W$.}
\item[\cercle{b}] \textit{For $Y$, $Y_i$ as in} \cercle{a}, \textit{for every}

\msp{page 125}

\textit{arrow $f\ :\ X\longrightarrow X'$ in $\Cat/Y$, we have $f\ \in\ W$ if and only if $\forall\ i\ \in\ I$, the morphism $f_i : X_i\longrightarrow X'_i$ over $Y_i$ is in $W$.}
\end{enumerate}

\textbf{Corollary 9.6.} \textit{If the axiom} Loc\,(5) \textit{is satisfied, then the connectedness axiom} Loc\,(8) \textit{is equivalent to this: for every arrow $f\ :\ X\longrightarrow Y$ in $W$, with $Y$ $0$-connected, $X$ is $0$-connected.}

\printedpage{68}
% head: A. Grothendieck : Dérivateurs, Ch. XV, Édition des Universités de Montpellier II et Paris VII
% foot: Version 5 novembre 2010

N.B. The axioms Loc\,(8) (connectedness) and Loc\,(5\,bis) (of decomposition) seem to be of a similar nature, both being linked to the decomposition of the objects of Cat into connected components. (Nevertheless, I think they are independent.) I presume that we will never really have to work with $W$ which do not satisfy these two axioms\,! (In addition to Loc\,(1, 2, 3), which for their part constitute the ``vital minimum''.) But apparently there is no need for them in the general theoretical developments so far. (But the connectedness axiom Loc\,(8) could play a role, when one speaks of the reciprocal determination of the structures of contractibility, and of $W$-asphericity, \textit{cf.} Pursuing Stacks. % [sic] "(Mais l'axiome ... Pursuing Stacks." -> closing parenthesis missing, the parenthesis opened at "(Mais" is never closed

\msp{page 126}

Let us note this:

\textbf{Proposition 9.7.}\footnote[50]{\textit{Cf.} stronger version in 9.9 below (rid of the hypothesis Loc\,(5).)} \textit{Suppose that $W$ satisfies} Loc\,(1, 2, 3, 5), i.e. \textit{it is a basic localizer satisfying the decomposition condition} (\textit{weak form}). \textit{If the axiom} Loc\,(8) \textit{of connectedness is} not \textit{satisfied, then the subcategory} $\mathrm{Hot}^*_W$ \textit{of the $X\ \in$} Cat \textit{which are nonempty is equivalent to the one-point category} (\textit{hence} $\mathrm{Hot}_W$ \textit{is equivalent either to the one-point category, if $\emptyset\longrightarrow e$ is an isomorphism in} $\mathrm{Hot}_W$, \textit{or to $\boldsymbol{\Delta}_1$ otherwise}\footnote[51]{The first case is also the one where $\exists\ f\in W$ with $\pi_0(f)$ \textit{not surjective,} or again $\emptyset\longrightarrow X$ in $W$, with $X$ nonempty.}.)

This is a variant of th. 2 of XII (p. 20), where we assumed Loc\,(6\,c$'$), stronger than Loc\,(5) and even than Loc\,(5\,bis), but with the stronger conclusion that $W=\Fl(\Cat)$.

\textsc{Proof}. Let
\[
f\ :\ X\longrightarrow Y\ ,\qquad f\ \in\ W\ ,
\]
be such that $\pi_0(f)$ is not \textit{surjective.} If $Y_0$ is a connected component which is not in $\mathrm{Im}\pi_0(f)$, then $X_0=X|Y_0$ is empty, and by Loc\,(7), $\emptyset\longrightarrow Y_0$ is in $W$. Thus $\forall\ Z$ in Cat, % [sic] Loc (7) -> Loc (5), apparently: the step X_0 -> Y_0 in W uses the decomposition hypothesis Loc (5) of 9.7, not Loc (7)
$\underbrace{\emptyset\times Z}_{=\,\emptyset}\longrightarrow Y_0\times Z$ is in $W$, and for every arrow $g : Z\longrightarrow Z'$, $\id_{Y_0}\times g : Y_0\times Z\longrightarrow Y_0\times Z'$ is in $W$. Taking $g$ such that $\pi_0(g)$ is not injective, e.g. $e\amalg e\longrightarrow e$, we find an arrow $\varphi=\id_{Y_0}\times g$ in $W$ such that $\pi_0(\varphi)$ is not injective.

\msp{page 127}

Let then $\varphi\ :\ I\longrightarrow J$ be in $W$, with $\pi_0(\varphi)$ not injective. Proceeding as before using Loc\,(5), we are reduced to the case where in addition $J$ is $0$-connected. (\textbf{N.B.} In both cases, we could also have invoked 9.6.) Let us choose two elements $i_0$, $i_1$ in $I$ in two distinct connected components, let $j_0$, $j_1$ be their images in $J$. We are going to treat $(I,i_0,i_1)$ and $(J,j_0,j_1)$ as ``homotopy intervals''. The fact that $i_0$, $i_1$ belong to distinct components of $I$ implies that two arrows $f_0,f_1 : X\rightrightarrows Y$ are ``homotopic with respect to $(I,i_0,i_1)$'', \textit{i.e.} there exists $h\ :\ X\times I\longrightarrow Y$ in Cat such that $f_0\ =\ h\partial_0^X$, $f_1=h\partial_1^X$, where
\[
\partial_0^X\ ,\ \partial_1^X\ :\ X\ \longrightarrow\ I\times X
\]
are given by $x\longmapsto(i_\alpha,x)$ $(\alpha\in\{0,1\})$. To find $h$, we decompose $I$ as
\[
I\ =\ I_0\amalg I_1\ ,\qquad\quad \text{with}\quad i_0\in I_0,\ i_1\in I_1\ ,
\]
whence $I\times X\ \simeq\ (I_0\times X)\amalg(I_1\times X)$, and we define $h$ by its restrictions $h_0$, $h_1$ to the two pieces $I_0\times X$, $I_1\times X$, namely
\[
h_0\ =\ f_0\circ\mathrm{pr}_2^{I_0}\ ,\qquad\quad h_1\ =\ f_1\circ\mathrm{pr}_2^{I_1}\ ,
\]

\end{document}
